\section*{Chapitre \ref{ch:g10:synthesis-yield} --- La synthèse~: rendement et pureté}

\begin{solution}{exo:g10:synthesis-yield:1}
$\eta = 4.2 / 5.0 = 0.84$, soit un \omterm{def:g10:synthesis-yield:yield}{rendement} de \qty{84}{\percent}.
\end{solution}

\begin{solution}{exo:g10:synthesis-yield:2}
Pesée des \omterm{def:g7:chemical-reactions:reactant}{réactifs}~; \omterm{def:g10:synthesis-yield:reflux}{chauffage à reflux}~; \omterm{def:g10:synthesis-yield:vacuum-filtration}{filtration sous vide} du
\omterm{def:g10:synthesis-yield:synthesis}{produit brut}~; \omterm{def:g10:synthesis-yield:recrystallisation}{recristallisation}~; mesure de la température de fusion.
\end{solution}

\begin{solution}{exo:g10:synthesis-yield:3}
Alimentée par le bas, l'eau remplit toute la double enveloppe avant de
sortir par le haut~: toute la longueur du réfrigérant est refroidie. Le
haut reste ouvert pour que le montage ne soit pas fermé~: chauffer un
montage fermé ferait monter la pression et pourrait le faire éclater.
\end{solution}

\begin{solution}{exo:g10:synthesis-yield:4}
Le produit doit être très soluble dans le \omterm{def:g6:solutions-and-solubility:solute}{solvant} chaud et très peu
soluble dans le \omterm{def:g6:solutions-and-solubility:solute}{solvant} froid. Les impuretés, présentes en petites
quantités, restent dissoutes dans le \omterm{def:g6:solutions-and-solubility:solute}{solvant} froid et finissent dans le
\omterm{def:g3:separating-mixtures:filter}{filtrat}.
\end{solution}

\begin{solution}{exo:g10:synthesis-yield:5}
Non~: il fond au-dessous de \qty{135}{\celsius} et sur un intervalle de
plusieurs degrés, signe d'un solide impur. Il faut le recristalliser, le
sécher et mesurer de nouveau sa température de fusion.
\end{solution}

\begin{solution}{exo:g10:synthesis-yield:6}
$n = 2.00 / 138.0 = \qty{0.0145}{mol}$ d'acide salicylique, le \omterm{def:g10:reaction-progress-table:limiting}{réactif
limitant}~; $m_{\max} = 0.0145 \times 180.0 = \qty{2.61}{g}$~;
$\eta = 2.03 / 2.61 = 0.78$, soit \qty{78}{\percent}.
\end{solution}

\begin{solution}{exo:g10:synthesis-yield:7}
La masse pesée comprend de l'eau~: elle est plus grande que la masse
d'aspirine, et le \omterm{def:g10:synthesis-yield:yield}{rendement} calculé est trop élevé, ici au-delà de
\qty{100}{\percent}, ce qui est impossible. Le \omterm{def:g10:synthesis-yield:synthesis}{produit brut} sec contient
encore des impuretés (acide salicylique qui n'a pas réagi,
sous-produits), qui ajoutent elles aussi de la masse~: son \omterm{def:g10:synthesis-yield:yield}{rendement} est
lui aussi trop élevé par rapport au vrai \omterm{def:g10:synthesis-yield:yield}{rendement} en aspirine pure.
\end{solution}

\begin{solution}{exo:g10:synthesis-yield:8}
Elle est beaucoup plus rapide, et elle laisse les cristaux bien plus
secs, puisque l'aspiration retire la plus grande partie du liquide~; les
cristaux sont aussi faciles à laver sur l'entonnoir et à racler sur le
papier filtre plat.
\end{solution}

\begin{solution}{exo:g10:synthesis-yield:9}
Produit~: $3.1/5.0 = 0.62$~; aspirine~: $3.1/5.0 = 0.62$~; acide
salicylique~: $2.2/5.0 = 0.44$. Le produit migre comme l'aspirine et ne
montre aucune tache d'acide salicylique~: c'est de l'aspirine, sans
acide salicylique détecté.
\end{solution}

\begin{solution}{exo:g10:synthesis-yield:10}
Le \omterm{def:g6:solutions-and-solubility:solute}{solvant} A~: très soluble à chaud, très peu soluble à froid~; le
\omterm{def:g6:solutions-and-solubility:solute}{solvant} B garde la plus grande partie du produit dissoute, même à froid.
Avec A, \qty{3.0}{g} demandent environ $3.0 / 150 = \qty{0.020}{L}$, soit
\qty{20}{mL} de \omterm{def:g6:solutions-and-solubility:solute}{solvant} chaud~; une fois froid, au plus
$2 \times 0.020 = \qty{0.04}{g}$ reste dissous.
\end{solution}

\begin{solution}{exo:g10:synthesis-yield:11}
Un liquide qui bout a besoin de place au-dessus de lui~: la mousse et
les projections ne doivent pas atteindre le réfrigérant. La pierre ponce
offre à la vapeur de petits sites où former des bulles, de sorte que le
liquide bout régulièrement au lieu de le faire par à-coups brusques.
\end{solution}

\begin{solution}{exo:g10:synthesis-yield:12}
$0.80 \times 0.70 = 0.56$~: \qty{56}{\percent}. Trois étapes à
\qty{90}{\percent}~: $0.90^3 = 0.73$, \qty{73}{\percent}. Chaque étape
fait perdre du produit (ainsi que du temps et de l'argent), et les pertes
se multiplient~: moins il y a d'étapes, plus il reste de produit à la
fin.
\end{solution}

\begin{solution}{exo:g10:synthesis-yield:13}
Au plus $3.3 \times 0.050 = \qty{0.17}{g}$ d'aspirine. L'eau glacée
dissout encore moins d'aspirine, car la \omterm{def:g6:solutions-and-solubility:solubility}{solubilité} diminue avec la
température.
\end{solution}

\begin{solution}{exo:g10:synthesis-yield:14}
\begin{enumerate}
  \item $n_0(\text{acide salicylique}) = 2.76 / 138.0 = \qty{0.0200}{mol}$,
        anhydride \qty{0.0500}{mol}~: l'acide salicylique est limitant,
        $x_{\max} = \qty{0.0200}{mol}$, et il reste $0.0500 - 0.0200 =
        \qty{0.0300}{mol}$ d'anhydride.
  \item La \omterm{def:g10:synthesis-yield:synthesis}{synthèse} donne \qty{0.0200}{mol} d'acide éthanoïque, la
        destruction de l'excès $2 \times 0.0300 = \qty{0.0600}{mol}$~:
        \qty{0.0800}{mol} en tout.
\end{enumerate}
\end{solution}

\begin{solution}{exo:g10:synthesis-yield:15}
$n(\text{aspirine}) = \num{1.00e6} / 180.0 = \qty{5.56e3}{mol}$. Avec un
\omterm{def:g10:synthesis-yield:yield}{rendement} de \qty{85}{\percent}, il faut en acide salicylique
$\num{5.56e3} / 0.85 = \qty{6.54e3}{mol}$, soit
$\num{6.54e3} \times 138.0 = \qty{9.0e5}{g}$, environ \qty{0.90}{t}.
\end{solution}

\begin{solution}{pb:g10:synthesis-yield:1}
\textbf{1.} C~: $7 + 4 = 11$ et $9 + 2 = 11$~; H~: $6 + 6 = 12$ et
$8 + 4 = 12$~; O~: $3 + 3 = 6$ et $4 + 2 = 6$. L'équation est ajustée.

\textbf{2.} Acide salicylique~: $7 \times 12.0 + 6 \times 1.0 + 3 \times 16.0
= \qty{138.0}{g/mol}$~; anhydride~: $4 \times 12.0 + 6 \times 1.0 +
3 \times 16.0 = \qty{102.0}{g/mol}$~; aspirine~: $9 \times 12.0 + 8 \times
1.0 + 4 \times 16.0 = \qty{180.0}{g/mol}$~; acide éthanoïque~: $2 \times 12.0
+ 4 \times 1.0 + 2 \times 16.0 = \qty{60.0}{g/mol}$.

\textbf{3.} \omterm{def:g7:chemical-reactions:reactant}{Réactifs}~: l'acide salicylique et l'anhydride éthanoïque~;
produit recherché~: l'aspirine~; sous-produit~: l'acide éthanoïque.

\textbf{4.} Chauffer accélère la réaction~; à reflux, les vapeurs sont
condensées et renvoyées dans le ballon, si bien qu'aucun \omterm{def:g7:chemical-reactions:reactant}{réactif} ni
produit n'est perdu.

\textbf{5.} $n = 3.0 / 138.0 = \qty{0.0217}{mol}$.

\textbf{6.} $m = 6.0 \times 1.08 = \qty{6.48}{g}$~;
$n = 6.48 / 102.0 = \qty{0.0635}{mol}$.

\textbf{7.} Acide salicylique~: $0.0217 - x$~; anhydride~: $0.0635 - x$~;
aspirine et acide éthanoïque~: $x$. L'acide salicylique est épuisé le
premier~: $x_{\max} = \qty{0.0217}{mol}$, l'acide salicylique est
limitant.

\textbf{8.} $m_{\max} = 0.0217 \times 180.0 = \qty{3.91}{g}$.

\textbf{9.} Pour que tout l'acide salicylique soit transformé en
aspirine~: s'il en restait, il resterait mêlé à l'aspirine et serait
difficile à éliminer, alors que l'excès d'anhydride est détruit par l'eau
ajoutée à la fin, et que l'acide éthanoïque qu'il donne est éliminé au
lavage.

\textbf{10.} Elle comprend de l'eau~: \qty{3.6}{g} donneraient un faux
\omterm{def:g10:synthesis-yield:yield}{rendement} de $3.6 / 3.91 = \qty{92}{\percent}$.

\textbf{11.} $3.3 / 3.91 = 0.84$~: \qty{84}{\percent} en \omterm{def:g10:synthesis-yield:synthesis}{produit brut} sec,
qui n'est pas pur.

\textbf{12.} Les impuretés sont éliminées, et une partie de l'aspirine
reste dissoute dans le \omterm{def:g6:solutions-and-solubility:solute}{solvant} froid et dans les eaux de lavage.

\textbf{13.} $3.3 \times 0.030 = \qty{0.10}{g}$ au plus.

\textbf{14.} $\eta = 2.7 / 3.91 = 0.69$~: \qty{69}{\percent}.

\textbf{15.} Une fusion nette à la valeur de l'aspirine pure~: les
cristaux sont de l'aspirine, et ils sont purs.

\textbf{16.} Une fusion au-dessous de \qty{135}{\celsius} et sur un large
intervalle~: le \omterm{def:g10:synthesis-yield:synthesis}{produit brut} était impur.

\textbf{17.} Le \omterm{def:g10:synthesis-yield:synthesis}{produit brut} contenait de l'aspirine et de l'acide
salicylique qui n'avait pas réagi~; la \omterm{def:g10:synthesis-yield:recrystallisation}{recristallisation} a éliminé
l'acide salicylique.

\textbf{18.} $R_f = 3.1 / 5.0 = 0.62$.

\textbf{19.} À peu près autant dans chacune~: de \qty{3.91}{g} possibles
à \qty{3.3}{g} de \omterm{def:g10:synthesis-yield:synthesis}{produit brut}, et de \qty{3.3}{g} à \qty{2.7}{g} lors
de la \omterm{def:g10:synthesis-yield:recrystallisation}{recristallisation}, soit environ \qty{0.6}{g} chaque fois. Comme le
\omterm{def:g10:synthesis-yield:synthesis}{produit brut} n'était pas de l'aspirine pure, la première partie du
travail a en réalité fait perdre un peu plus de \qty{0.6}{g}
d'aspirine.

\textbf{20.} Le \omterm{def:g10:synthesis-yield:yield}{rendement} de la \omterm{def:g10:synthesis-yield:synthesis}{synthèse} est d'environ
\qty{69}{\percent}.
\end{solution}
